% c11.fig.drehung — the rotated coordinate of step 2 in the proof of Proposition c11.prop.zweikoerper.
%
% Generating formula (REVIEW.md A2; chapter equation c11.eq.drehung), for one Cartesian component:
%     z = (x_e/sqrt(2 m_e) - x_p/sqrt(2 m_p))/gamma,   gamma = sqrt(1/(2 m_e) + 1/(2 m_p)),
% that is  z = sqrt(m_p/M) x_e - sqrt(m_e/M) x_p,  a unit vector (sqrt(m_p/M), -sqrt(m_e/M)) in the
% plane of the scaled coordinates (x_e, x_p).  Drawn with the STATED, NOT PHYSICAL ratio
% m_p : m_e = 3 : 1:  direction of z = (sqrt 3/2, -1/2), i.e. the x_e-axis turned by -30 degrees;
% the perpendicular direction is (1/2, sqrt 3/2).  Both new axes and the unit circle are drawn.
% Greyscale: old axes thin, new axes heavy, labels.
\begin{tikzpicture}[x=1.7cm, y=1.7cm, >={Stealth[length=5pt,width=4pt]}]
  \draw[tpc axis, ->] (-1.5,0) -- (1.7,0) node[right, tpc label] {$x_{\mathrm{e}}$};
  \draw[tpc axis, ->] (0,-1.25) -- (0,1.45) node[above, tpc label] {$x_{\mathrm{p}}$};
  \draw[tpc leader, dash pattern=on 2pt off 2pt] (0,0) circle (1);
  \draw[draw=tpc-plate-ink, line width=\tpcstrokestate, ->] ({-1.4*cos(30)},{1.4*sin(30)}) -- ({1.75*cos(30)},{-1.75*sin(30)})
    node[right, tpc label] {$z$};
  \draw[draw=tpc-plate-ink, line width=\tpcstrokestate, ->] ({-1.2*sin(30)},{-1.2*cos(30)}) -- ({1.5*sin(30)},{1.5*cos(30)});
  \fill[tpc-plate-ink] ({cos(30)},{-sin(30)}) circle (1.6pt);
  \node[tpc label, anchor=west] at (1.12,-0.2)
    {$\bigl(\sqrt{m_{\mathrm{p}}/M},\,-\sqrt{m_{\mathrm{e}}/M}\bigr)$};
  \draw[tpc leader] (1.12,-0.22) -- ({cos(30)+0.03},{-sin(30)+0.03});
  \draw[tpc leader, ->] (0.55,0) arc[start angle=0, end angle=-30, radius=0.55];
  \node[tpc label, anchor=south west, align=left, text=tpc-plate-muted] at (1.0,0.75)
    {\tpcifenglish{drawn for}{gezeichnet für}\\ $m_{\mathrm{p}} : m_{\mathrm{e}} = 3 : 1$};
\end{tikzpicture}
